% Generated draft: requires mathematical and source review.
\documentclass[11pt]{article}
\usepackage[margin=25mm]{geometry}
\usepackage{fontspec}
\setmainfont{Latin Modern Roman}
\usepackage{amsmath,amssymb,mathtools,unicode-math}
\setmathfont{Latin Modern Math}
\usepackage{enumitem,needspace,xurl,hyperref}
\setlength{\parindent}{0pt}
\setlength{\parskip}{5pt}
\setlength{\emergencystretch}{3em}
\newcommand{\N}{\mathbb{N}}
\newcommand{\Z}{\mathbb{Z}}
\newcommand{\Q}{\mathbb{Q}}
\newcommand{\R}{\mathbb{R}}
\newcommand{\C}{\mathbb{C}}
\newcommand{\F}{\mathbb{F}}
\newcommand{\A}{\mathbb{A}}
\newcommand{\PP}{\mathbb{P}}
\newcommand{\E}{\mathbb{E}}
\newcommand{\eps}{\varepsilon}
\newcommand{\dd}{\,\mathrm{d}}
\begin{document}
\section*{UnsolvedMath ID: 800329}
\textbf{Source problem code:} catalogue entry
\subsection{Definitions and mathematical statement}
Let \(i\in\C\) satisfy \(i^2=-1\). For every integer \(n\geq 0\), define the \emph{Gessel sequence} by
\[
b_n=2\cdot 5^n-(3+4i)^n-(3-4i)^n.
\]
The index range is supported by the source's displayed generating series, which starts at \(n=0\). The source states that \(b_n\in\Z_{\geq 0}\) for every \(n\geq 0\), and that the sequence is C-recursive, meaning that it satisfies a linear recurrence with constant coefficients. It presents this sequence as an example without an easy combinatorial interpretation.

\textbf{Unresolved conjectural statement.} The supplied extraction of Conjecture 13 ends with a broken sentence referring to a sequence and to type \(\mathrm{(C1)}\). It supplies neither the definition of this type nor a complete assertion about membership in it. Consequently, the sequence is precisely specified above, but the intended conjecture cannot be reconstructed as a definite mathematical statement from the supplied record. The displayed integrality expansion and generating-function formula are also truncated and are not reproduced as valid identities.
\subsection{Short English statement}
The record defines the Gessel sequence \(b_n=2\cdot5^n-(3+4i)^n-(3-4i)^n\) for \(n\geq0\). Its intended conjecture concerns type \(\mathrm{(C1)}\), but the extraction omits both the type's definition and the complete conjectural assertion. This statement-only draft therefore requires source recovery and human mathematical review.
\subsection{Sources}
\begin{itemize}
\item Local UnsolvedMath registry, permanent ID \(800329\), source problem code \texttt{catalogue entry}; collection: AMR Open Problem Lists. The supplied record attributes the problem to Pak, with proposed year 2018, and labels it \texttt{open}. This status is supplied registry metadata, not an independently verified literature assessment.
\item Pak, \emph{Complexity Problems in Enumerative Combinatorics} (2018), Conjecture 13, as identified in the supplied source list: \url{https://arxiv.org/abs/1803.06636}. The record reports extraction by \texttt{arxiv\_latex} and access on 2026-10-16; these are supplied provenance details, not a claim of independent source access or verification.
\item The source gives the attribution/context ``cf. [Ges-sem]'' and refers to that item for the definition of the Gessel sequence. No bibliographic identification or URL for [Ges-sem] is supplied.
\item OEIS entry A250102, linked in the supplied statement: \url{https://oeis.org/A250102}. Its contents have not been independently checked for this draft.
\end{itemize}
\end{document}
